\documentclass[10pt,a4paper]{amsart}
\usepackage[margin=19mm]{geometry}
\usepackage{amsmath,amssymb,mathtools}
\usepackage[scr=rsfs,cal=euler]{mathalfa}
\usepackage{xfrac}
\usepackage[unicode,hidelinks,bookmarks=false]{hyperref}
\newcommand{\ie}{i.e.\ }
\newcommand{\cf}{cf.\ }
\newcommand{\resp}{respectively\ }
\newcommand{\Subsection}{\subsection}
\providecommand{\nobreakdash}{\nobreak\hbox{-}}
\setlength{\emergencystretch}{2em}
\multlinegap0pt
\usepackage{bm}
\usepackage{csquotes}
\usepackage{amssymb}
\usepackage{booktabs}
\usepackage{multirow}
\usepackage{caption}
\usepackage{algorithm}\usepackage{algpseudocode}
\usepackage{xcolor}
\newtheorem{theorem}{Theorem}
\newcommand{\E}{\mathbb{E}}
\newcommand{\G}{\widehat{\mathbf{\Gamma}}}
\newcommand{\N}{\mathcal{N}}
\newcommand{\Si}{\mathbf{\Sigma}}
\title{\textit{Carr\'{e} du champ} flow matching: better quality-generalisation tradeoff in generative models}
\author{Jacob Bamberger, Iolo Jones, Dennis Duncan, Michael M. Bronstein, Pierre Vandergheynst, Adam Gosztolai}
\date{Source: arXiv:2510.05930v2 (CC BY 4.0)}
\begin{document}
\maketitle

\noindent\emph{Source and licence.} \textit{Carr\'{e} du champ} flow matching: better quality-generalisation tradeoff in generative models. Version arXiv:2510.05930v2 (CC BY 4.0), distributed by arXiv under the Creative Commons Attribution 4.0 International licence, http://creativecommons.org/licenses/by/4.0/, as recorded in the arXiv licence metadata for this exact version and shown on the abstract page https://arxiv.org/abs/2510.05930v2. Full licence notice: \texttt{licenses/arxiv-2510.05930-CC-BY-4.0.txt}.\par\smallskip

\noindent\textit{Editorial scope.} Counted unit: the article's theorem optimalGM (source0.tex lines 1414--1429). The article's complete original proof of it is delivered with it (source0.tex lines 1431--1458, one \texttt{proof} environment of 28 lines). No mathematics is written, repaired or re-derived: the statement and the proof are the article's own lines, in the article's own notation, and no retained line is substituted or re-flowed.  Retained uncounted as the source-local context the proof needs: lines 785--788.  Nothing was omitted from any retained span, so the package carries no omission record.  No retained line contains a citation, so the wrapper carries no bibliography and no bibliography entry.  No retained line contains a figure environment, an image-inclusion command or drawing code, so no image is delivered and the package declares no asset dependency.  Editorial formatting only, in addition: an AMS \texttt{amsart} preamble that restates the article's own declarations and macros used by the retained lines, namely the environments \texttt{theorem} on separate counters, together with the standard packages those lines need.  The retained environments print in this document as Theorem 1 in order of appearance, whereas the article prints the same items with the numbering of its own full document.  The complete native source of the article as distributed in the arXiv e-print for this exact version is bundled unchanged as \texttt{sources/186570/source0.tex}.
\begin{equation}\label{eq:markov_kernel}
    (Pf)(x^{(i)}) := \sum_j P_{ij}\,f(x^{(j)}) \;=\; \E_{Y\sim P_i}\!\big[f(Y)\big],
\qquad P_{ij} := \frac{w_\epsilon(x^{(i)},x^{(j)})}{\sum_\ell w_\epsilon(x^{(i)},x^{(\ell)})}.
\end{equation}

\begin{theorem} \label{optimalGM}
Let $P_{xy}$ be a transition kernel on $\mathbb{R}^d$. Let $P_x$ denote the probability measure supported on the $k$ nearest neighbours of $x$ obtained by restricting and renormalising $P_{xy}$ (\ref{eq:markov_kernel}).
Assume $P_x$ has finite second moments.
Then the unique maximiser over $m\in\mathbb{R}^d$ and $\Sigma$ over the space of positive definite matrices of the expected log-likelihood of a sample $Y$ from a Gaussian density \( \N(m, \Si) \)
\[
(m,\Si) \;\longmapsto\; \E_{Y\sim P_x}\!\left[\log \mathcal{N}(Y; m,\Sigma)\right]
\]
is given by matching the first two moments of \( P_x \):
\[
m^*(x)=\E_{Y\sim P_x}[Y],
\qquad
\Si^*(x):= \G(x) = \E_{P_x}\!\left[(Y-m^*)(Y-m^*)^\top\right].
\]
Equivalently, $\mathcal{N}(m^*,\Si^*)$ minimises $\mathrm{KL}(P_x\,\Vert\,\mathcal{N}(m,\Si))$.

\end{theorem}

\begin{proof}
For $\Sigma\succ0$, the log-likelihood of a sample $Y\sim \mathcal{N}(m, \Si)$ is
\[
\log \N(Y) = -\frac{d}{2}\log(2\pi) - \frac{1}{2}\log\det \Si - \frac{1}{2}(Y - m)^\top \Si^{-1} (Y - m).
\]
Taking the expectation with respect to $Y \sim P_x$ gives
\[
\E_{Y \sim P_x}\left[\log \N(Y; m, \Si)\right] = -\frac{d}{2}\log(2\pi) - \frac{1}{2}\log \det\Si - \frac{1}{2}\E_{Y \sim P_x}\left[(Y - m)^\top \Si^{-1} (Y - m)\right].
\]
We can differentiate with respect to $m$ to get
\[
\frac{\partial}{\partial m} \E_{Y \sim P_x}\left[\log \N(Y; m, \Si)\right] = \Si^{-1} (\E_{Y \sim P_x}[Y] - m),
\]
which is zero at the maximum $m^*(x) = \E_{Y \sim P_x}[Y]$.
To find the optimal $\Si$, let $S_x=\E[(Y-m^*)(Y-m^*)^\top]$. We can rewrite the expectation term as 
\[
\E[\log \mathcal{N}(Y;m^*,\Si)]
= \tfrac{1}{2}\log\det\Si^{-1} - \tfrac{1}{2}\mathrm{tr}(\Si^{-1} S_x) + \text{const},
\]
which is strictly concave in $\Si^{-1}\succ0$.
% $\text{tr}(\Si^{-1} \E_{Y \sim P_x}[(Y - m^*)(Y - m^*)^\top])$.
Differentiating with respect to $\Si^{-1}$ we obtain
\[
\frac{\partial}{\partial \Si^{-1}} \E_{Y \sim P_x}\left[\log \N(Y; m^*, \Si)\right] = \frac{1}{2}\Si - \frac{1}{2}S_x.
\]
This quantity is zero at the maximum $\Si^* = \E_{Y \sim P_x}\left[(Y - m^*)(Y - m^*)^\top\right]$, and uniqueness follows from strict concavity. 
Thus, the optimal covariance is given by $\G := \E_{Y \sim P_x}\left[(Y - m^*)(Y - m^*)^\top\right]$, the centred covariance of the neighbours of $x$.
\end{proof}

\end{document}
